\scanpage{104}
% archive foliation: 103
\deleted{
(par le choix d'isomorphismes $\unsure{\widehat{\mathbf Z}_\ell|\mathbf Z_\ell}(1)\simeq\mathbf Z_\ell$ pour $\forall\ell$)

\unsure{on trouve} qu'on a défini une représentation du produit semi-direct de $\mathbf Z$ avec $\mathbf Z$, opérant sur le sous-groupe $\mathbf Z$ par multiplication par $q$. Donc ce groupe est engendré par deux éléments $f$ (correspondant à $\operatorname{frob}_q$) et $g$, satisfaisant la relation
\[
(6.4)\qquad fgf^{-1}=g^q
\]
\unsure{Or on a le}

\uline{Lemme} 6.5. Soient $F,G$ deux automorphismes d'un vectoriel \inserted{$E$} de dim finie sur un corps, et $q$ \deleted{supposant que} un entier $>1$ tel qu'on ait \deleted{\unsure{$FG$}}
\[
FGF^{-1}=G^q
\]
Alors $G$ est quasi-unipotent, i.e. ses valeurs propres sont toutes des racines de $1$, ou encore: il existe un entier $N\geq1$ tel que $G^N$ soit unipotent.

\uline{Dém.} L'\unsure{ens} des valeurs propres de $G$ \deleted{et $FGF^{-1}$} \deleted{dans} (\unsure{c-à-d vues} dans une clôture alg. $\overline R$ de $R$) étant le \deleted{\unsure{même, et}} $=$ que celui de $FGF^{-1}$, donc le $=$ que celui de $G^q$, il s'ensuit que c'est une partie \inserted{\deleted{\unsure{finie}} $\Phi$} de $\overline R^*$ stable par l'application
\[
\lambda\longrightarrow\lambda^q
\]
Comme $\Phi$ est fini, il existe donc pour tt $\lambda\in\Phi$ deux entiers $N'\unsure{>\!|\geq}N''$ \deleted{tels} $\geq0$ tels que
}
